% Excerpts from Marcolli-Manin formalization test document
% Based on themes from "Continued Fractions, Modular Symbols, and Noncommutative Geometry"

\documentclass{article}
\usepackage{amsmath,amssymb}

\title{Modular Symbols and Noncommutative Spaces}
\author{Test formalization document}

\begin{document}

\section{Introduction}
\label{sec:intro}

This document contains excerpts for testing the autoformalizer protocol
with mathematical content inspired by the work of Marcolli and Manin on
modular symbols and noncommutative geometry.

\section{Modular Curves}
\label{sec:modular-curves}

Let $\mathbb{H}$ denote the upper half-plane $\{z \in \mathbb{C} : \Im(z) > 0\}$.
The modular group $\text{SL}_2(\mathbb{Z})$ acts on $\mathbb{H}$ by M\"obius transformations:

$$\begin{pmatrix} a & b \\ c & d \end{pmatrix} \cdot z = \frac{az + b}{cz + d}$$

\subsection{Fundamental Domain}
\label{sec:fundamental-domain}

\begin{definition}[Standard Fundamental Domain]
\label{def:fundamental-domain}
The standard fundamental domain $\mathcal{F}$ for the action of $\text{SL}_2(\mathbb{Z})$
on $\mathbb{H}$ is:
$$\mathcal{F} = \{z \in \mathbb{H} : |z| \geq 1, |\Re(z)| \leq \tfrac{1}{2}\}$$
\end{definition}

\begin{theorem}[Fundamental Domain Property]
\label{thm:fundamental-domain}
Every point $z \in \mathbb{H}$ is equivalent under $\text{SL}_2(\mathbb{Z})$ to a unique
point in the closure of $\mathcal{F}$, up to identifications on the boundary.
\end{theorem}

\begin{proof}
The proof proceeds by showing that for any $z \in \mathbb{H}$, there exists
$\gamma \in \text{SL}_2(\mathbb{Z})$ such that $\gamma \cdot z \in \overline{\mathcal{F}}$.
This uses the generators $S = \begin{pmatrix} 0 & -1 \\ 1 & 0 \end{pmatrix}$
and $T = \begin{pmatrix} 1 & 1 \\ 0 & 1 \end{pmatrix}$.
\end{proof}

\section{Modular Symbols}
\label{sec:modular-symbols}

\begin{definition}[Modular Symbol]
\label{def:modular-symbol}
A modular symbol is a formal sum of geodesics in $\mathbb{H}^* = \mathbb{H} \cup \mathbb{P}^1(\mathbb{Q})$
connecting cusps. For cusps $\alpha, \beta \in \mathbb{P}^1(\mathbb{Q})$, we write
$\{\alpha, \beta\}$ for the geodesic from $\alpha$ to $\beta$.
\end{definition}

\begin{theorem}[Three-term Relation]
\label{thm:three-term}
For any three cusps $\alpha, \beta, \gamma \in \mathbb{P}^1(\mathbb{Q})$:
$$\{\alpha, \beta\} + \{\beta, \gamma\} = \{\alpha, \gamma\}$$
\end{theorem}

\begin{proof}
This follows from the definition of modular symbols as formal sums of paths
and the concatenation of geodesics.
\end{proof}

\begin{definition}[Space of Modular Symbols]
\label{def:symbol-space}
Let $\mathcal{M}$ denote the $\mathbb{Q}$-vector space generated by symbols
$\{\alpha, \beta\}$ for $\alpha, \beta \in \mathbb{P}^1(\mathbb{Q})$, modulo:
\begin{enumerate}
\item $\{\alpha, \alpha\} = 0$
\item $\{\alpha, \beta\} + \{\beta, \gamma\} + \{\gamma, \alpha\} = 0$
\end{enumerate}
\end{definition}

\section{Continued Fractions}
\label{sec:continued-fractions}

\begin{definition}[Continued Fraction Expansion]
\label{def:cf-expansion}
For $x \in (0,1)$, the continued fraction expansion is:
$$x = \cfrac{1}{a_1 + \cfrac{1}{a_2 + \cfrac{1}{a_3 + \cdots}}}$$
where $a_i \in \mathbb{Z}_{> 0}$.
\end{definition}

\begin{theorem}[Gauss Map]
\label{thm:gauss-map}
The Gauss map $G : (0,1) \to (0,1)$ defined by $G(x) = \{1/x\}$ (fractional part)
generates the continued fraction digits: $a_n(x) = \lfloor 1/G^{n-1}(x) \rfloor$.
\end{theorem}

\begin{proposition}[Convergents]
\label{prop:convergents}
The convergents $p_n/q_n$ of a continued fraction satisfy:
\begin{align}
p_n &= a_n p_{n-1} + p_{n-2} \\
q_n &= a_n q_{n-1} + q_{n-2}
\end{align}
with $p_{-1} = 1, p_0 = 0, q_{-1} = 0, q_0 = 1$.
\end{proposition}

\section{Noncommutative Boundary}
\label{sec:nc-boundary}

\begin{definition}[Noncommutative Torus]
\label{def:nc-torus}
The noncommutative torus $\mathbb{T}^2_\theta$ is the universal C*-algebra
generated by unitaries $U, V$ satisfying:
$$VU = e^{2\pi i \theta} UV$$
\end{definition}

\begin{theorem}[Morita Equivalence]
\label{thm:morita}
For $\theta, \theta'$ related by an element of $\text{GL}_2(\mathbb{Z})$,
the noncommutative tori $\mathbb{T}^2_\theta$ and $\mathbb{T}^2_{\theta'}$
are Morita equivalent.
\end{theorem}

\section{Arithmetic and Dynamics}
\label{sec:arithmetic}

\begin{theorem}[Geodesic Coding]
\label{thm:geodesic-coding}
There is a bijection between oriented geodesics on the modular surface
$\text{SL}_2(\mathbb{Z}) \backslash \mathbb{H}$ and bi-infinite sequences
of positive integers $(a_n)_{n \in \mathbb{Z}}$.
\end{theorem}

\begin{corollary}[Periodic Geodesics]
\label{cor:periodic-geodesics}
Closed geodesics on the modular surface correspond to eventually periodic
continued fractions, i.e., quadratic irrationals.
\end{corollary}

\end{document}
